\section{Related Work}

\subsection{Localized Programming Languages and Non-English Syntax}

Prior work on localized programming languages and non-English syntax provides
the closest comparison point for \plat. Relevant examples include languages,
teaching environments, or dialects that replace English keywords, expose
localized commands, or support programming in scripts other than Latin.

% TODO: Replace \citep{todo-localized-programming-languages} with precise
% citations to localized programming languages and non-English syntax projects.
\citep{todo-localized-programming-languages}

\subsection{Educational and Beginner-Oriented Languages}

Beginner-oriented environments often simplify syntax, reduce incidental
complexity, and foreground immediate feedback. \plat{} shares those concerns
through its small feature set, direct script execution, and readable examples,
but its central focus is cultural localization rather than pedagogy alone.

% TODO: Add sources on Logo, Scratch, BASIC-family teaching languages, and
% beginner-oriented language design.
\citep{todo-educational-languages}

\subsection{Esoteric Languages and Cultural Expression}

Esolangs demonstrate that programming-language design can be artistic,
humorous, critical, or culturally expressive. \plat{} is not primarily an
esolang, but it overlaps with esolang practice insofar as it treats syntax and
vocabulary as expressive materials rather than purely ergonomic mechanisms.

% TODO: Add sources on esolangs, code art, and cultural expression in language
% design.
\citep{todo-esolangs-cultural-expression}

\subsection{Programming Languages as Sociotechnical Systems}

Programming languages are embedded in documentation, tooling, communities,
education, governance, and economic structures. This sociotechnical view helps
explain why localized keywords alone are insufficient: diagnostics, examples,
community feedback, and maintenance practices are part of the language system.

% TODO: Add sources from software studies, HCI, CSCW, and PL community studies.
\citep{todo-pl-sociotechnical}

\subsection{Digital Language Preservation}

Digital language preservation work asks how minority and regional languages can
remain visible and usable in contemporary computing environments. \plat{} is a
small programming-language artifact rather than a preservation platform, but
it participates in the broader question of whether software systems can create
meaningful domains of use for regional languages.

% TODO: Add sources on digital language preservation, minority-language
% technology, and language revitalization.
\citep{todo-digital-language-preservation}
